\documentclass[11pt,a4paper]{article}

\usepackage[T1]{fontenc}
\usepackage[utf8]{inputenc}
\usepackage[polish]{babel}
\usepackage{lmodern}
\usepackage[margin=2.5cm]{geometry}
\usepackage{booktabs}
\usepackage{float}
\usepackage{tabularx}
\usepackage[colorlinks, linkcolor=black, urlcolor=blue]{hyperref}

\newcolumntype{L}{>{\raggedright\arraybackslash}X}
\newcolumntype{R}{>{\raggedleft\arraybackslash}X}

\title{Porównanie parametrów technicznych samochodów\\[4pt]
\large Kamień milowy I}
\author{inż. Filip Stępień \and inż. Bartłomiej Karkoszka}
\date{Politechnika Świętokrzyska\\
Wydział Elektrotechniki, Automatyki i Informatyki\\
Przedmiot: Statystyka w informatyce\\[8pt]
listopad 2026}

\begin{document}

\maketitle

\section{Cel i zakres projektu}

Celem projektu jest interaktywna aplikacja do analizy statystycznej i wizualizacji danych
o samochodach. Pozwala ona porównywać pojazdy pod względem osiągów, kosztów eksploatacji,
funkcjonalności i ceny. Obejmuje miary rozkładu, korelacje, histogramy, wykresy punktowe,
wykresy obrazkowe oraz ranking aut według wag ustalanych przez użytkownika.

Na etapie kamienia milowego aplikacja realizuje wszystkie wymagania projektu.
Sprawozdanie przedstawia jej obecny stan oraz wstępne wyniki dla pełnego zbioru danych.

\section{Rodzaj i źródło danych}

Dane pochodzą ze zbioru \emph{Qatar Cars} (P.~Musgrave, 2025, licencja CC~BY~4.0),
dostępnego pod adresem \url{https://github.com/profmusgrave/qatarcars}. Zbiór zebrano
w styczniu 2025~r., głównie z serwisu YallaMotors Qatar. Jeden rekord opisuje wariant
samochodu określony marką, modelem, wersją i rokiem modelowym.

\begin{table}[H]
\centering
\begin{tabular}{lll}
\toprule
Atrybut & Jednostka & Uwagi \\
\midrule
Moc & KM & \\
Spalanie & l/100 km & \\
Pojemność bagażnika & l & część wartości wątpliwa \\
Cena & PLN & przeliczona z rialów katarskich \\
Przyspieszenie 0--100 km/h & s & zamiast prędkości maksymalnej \\
\bottomrule
\end{tabular}
\caption{Analizowane atrybuty}
\end{table}

Surowy plik zawiera 105 samochodów. Pominięto 10 aut elektrycznych, bo źródło nie podaje
dla nich zużycia paliwa. Analizowany zbiór liczy 95 samochodów 27 marek (81 benzynowych
i 14 hybryd), więc wymaganie co najmniej 30 obiektów i 5 atrybutów jest spełnione
z zapasem. Ceny przeliczono na złote po średnim kursie NBP z 15.01.2025
(1,1354~PLN/QAR). Przeliczenie liniowe nie zmienia skośności, kurtozy, współczynnika
zmienności ani korelacji.

Żaden darmowy zbiór nie zawierał wszystkich pięciu atrybutów z tematu projektu naraz.
Prędkość maksymalną zastąpiono przyspieszeniem 0--100~km/h, które również opisuje osiągi auta.

\section{Technologie}

\begin{table}[H]
\centering
\begin{tabular}{ll}
\toprule
Technologia & Zastosowanie \\
\midrule
Python 3.12 / Streamlit & logika i interfejs użytkownika \\
pandas / SciPy & przetwarzanie danych i miary statystyczne \\
Plotly & interaktywne wykresy \\
SQLite / SQLAlchemy & przechowywanie danych i dostęp do nich \\
Docker / Task & uruchamianie aplikacji jednym poleceniem \\
\bottomrule
\end{tabular}
\caption{Wykorzystane technologie}
\end{table}

Pięć analizowanych atrybutów jest opisanych w jednym miejscu kodu (nazwa, jednostka,
kierunek ,,mniej znaczy lepiej'', dokładność wyświetlania). Na tej podstawie powstają
kolumny bazy, statystyki, ranking i wykresy, więc nazwy i jednostki nie są powielane.

\section{Przechowywanie danych}

Dane są przechowywane w relacyjnej bazie SQLite, w jednej tabeli samochodów. Tabela zawiera
identyfikator, markę, model, wersję, rok, pięć atrybutów liczbowych oraz adres i datę źródła.
Ograniczenia w bazie pilnują poprawności danych: wszystkie atrybuty muszą być dodatnie
(bagażnik nieujemny), a ten sam wariant auta może wystąpić tylko raz.

Przy pierwszym uruchomieniu dane są importowane z pliku CSV w jednej transakcji. Ponowne
uruchomienie nie tworzy duplikatów. Baza plikowa wystarcza przy niewielkim zbiorze, jednym
użytkowniku i danych tylko do odczytu, a nie wymaga osobnego serwera.

\section{Analiza statystyczna}

Dla każdego atrybutu aplikacja wyznacza miary zestawione w tabeli~\ref{tab:miary}.

\begin{table}[H]
\centering
\begin{tabularx}{\textwidth}{lL}
\toprule
Grupa miar & Miary \\
\midrule
Tendencja centralna & średnia, mediana, dominanta \\
Zróżnicowanie & minimum, maksimum, rozstęp, kwartyle $Q_1$ i $Q_3$, rozstęp
międzykwartylowy, wariancja, odchylenie standardowe, współczynnik zmienności \\
Asymetria & współczynnik skośności \\
Koncentracja & kurtoza nadmiarowa \\
\bottomrule
\end{tabularx}
\caption{Wyznaczane miary rozkładu}
\label{tab:miary}
\end{table}

Wariancja i odchylenie standardowe są liczone dla próby (dzielnik $n-1$). Skośność wymaga
co najmniej trzech obserwacji, a kurtoza co najmniej czterech. Rozkład każdego atrybutu
można obejrzeć na histogramie połączonym z wykresem pudełkowym.

\section{Korelacje pomiędzy atrybutami}

Aplikacja wyznacza macierz korelacji $5 \times 5$ metodą Pearsona (zależność liniowa) lub
Spearmana (zależność monotoniczna, na rangach). Macierz jest pokazana jako tabela i mapa
cieplna, a aplikacja wskazuje parę atrybutów o najsilniejszej zależności. Korelacja nie
oznacza związku przyczynowego, a wyniki zależą od wybranych marek i zakresu cen.

\section{Wizualizacja danych}

Zakładka \emph{Statystyki} pokazuje tabelę miar, histogram i wykres pudełkowy.
W zakładce \emph{Porównania} użytkownik wybiera dwa atrybuty do wykresu punktowego,
na którym marki są rozróżnione kolorami.

Nietypowym elementem jest ,,garaż obrazkowy''. Użytkownik wybiera do sześciu samochodów
i jeden atrybut, a wartości są przedstawione jako rzędy ikon samochodu. Jedna ikona
odpowiada 25~KM, 1~l/100~km, 100~l, 25\,000~PLN lub 1~s. Liczba ikon jest zaokrąglana
i ograniczona do 200, dlatego obok wyświetlana jest dokładna wartość.

\section{Ranking samochodów}

Użytkownik ustala wagi pięciu atrybutów w skali od 0 do 10. Każdy atrybut jest
normalizowany do przedziału 0--1, a dla ceny, spalania i czasu przyspieszenia skala jest
odwracana, bo mniejsza wartość jest lepsza. Wynik końcowy to średnia ważona przeskalowana
do 0--100 punktów. Ranking jest liczony dla aktualnie wybranych aut.

\section{Interfejs}

Aplikacja działa w przeglądarce i ma pięć zakładek: \emph{Samochody}, \emph{Statystyki},
\emph{Porównania}, \emph{Korelacje} i \emph{Ranking}. Panel boczny pozwala filtrować auta
według marki i maksymalnej ceny; filtry wpływają na wszystkie analizy. Gdy po filtrowaniu
zostaje mniej niż 30 aut, aplikacja wyświetla ostrzeżenie.

\section{Wstępne wyniki}

Tabela~\ref{tab:wyniki} przedstawia wybrane miary rozkładu dla pełnego zbioru
(95 samochodów).

\begin{table}[H]
\centering
\small
\begin{tabularx}{\textwidth}{lRRRRR}
\toprule
Miara & Moc [KM] & Spalanie [l/100 km] & Bagażnik [l] & Cena [tys.\ PLN] & Przysp. [s] \\
\midrule
Średnia & 297,6 & 8,74 & 426,2 & 874,8 & 7,28 \\
Mediana & 240,0 & 7,60 & 440,0 & 160,4 & 7,00 \\
Odch. standardowe & 247,8 & 3,51 & 259,6 & 4\,126,8 & 2,88 \\
Wsp. zmienności & 83\% & 40\% & 61\% & 472\% & 40\% \\
Skośność & 3,23 & 1,72 & 0,70 & 8,01 & 0,51 \\
Kurtoza nadmiarowa & 14,21 & 3,91 & 0,82 & 68,53 & $-0{,}28$ \\
\bottomrule
\end{tabularx}
\caption{Wybrane miary rozkładu}
\label{tab:wyniki}
\end{table}

\begin{itemize}
  \item Cena ma skrajnie prawostronnie skośny rozkład. Średnia jest ponad pięć razy wyższa
  od mediany, bo podnosi ją kilka supersamochodów (np. Bugatti Centodieci za
  37,5~mln~PLN). Typową cenę lepiej opisują mediana i rozstęp międzykwartylowy.
  \item Moc również jest wyraźnie prawostronnie skośna, z tego samego powodu.
  \item Bagażnik i przyspieszenie mają rozkłady zbliżone do symetrycznych.
  \item Najsilniejsza zależność występuje między mocą a ceną (Pearson 0,765,
  Spearman 0,930). Moc i czas przyspieszenia są silnie ujemnie skorelowane
  (Spearman $-0{,}852$), a moc i spalanie dodatnio (Pearson 0,742).
  \item Bagażnik praktycznie nie zależy od pozostałych atrybutów
  (Spearman od $-0{,}04$ do 0,14).
  \item Współczynnik Spearmana lepiej opisuje ten zbiór niż Pearsona, bo jest odporny
  na wartości skrajne cen i mocy.
\end{itemize}

\section{Uruchamianie}

Do uruchomienia potrzebne są
\href{https://www.docker.com/products/docker-desktop/}{Docker Desktop}
(z Docker Compose w wersji 2) oraz narzędzie \href{https://taskfile.dev}{Task}.

\begin{enumerate}
  \item Uruchomić Docker Desktop.
  \item Przejść do katalogu projektu.
  \item Wykonać polecenie \texttt{task dev}. Obraz zostanie zbudowany, a aplikacja uruchomiona.
  \item Otworzyć w przeglądarce adres \url{http://localhost:8501}.
\end{enumerate}

Przy pierwszym uruchomieniu baza danych jest tworzona automatycznie i wypełniana danymi
z pliku CSV. Aplikację zatrzymuje polecenie \texttt{task stop}; przy kolejnym starcie baza
powstaje od nowa z tego samego pliku.

\section{Źródła}

\begin{itemize}
  \item Repozytorium projektu: \url{https://github.com/codingbart/car-metrics}
  \item Zbiór Qatar Cars: \url{https://github.com/profmusgrave/qatarcars}
  \item Dokumentacja: Python, pandas, SciPy, Plotly, Streamlit, SQLAlchemy.
\end{itemize}

\end{document}
